\chapter{Differentials,  Differentiation, and the Derivative }
\label{chapt:differentials}
\mycenterquote{By relieving the brain of all unnecessary work, a good
  notation sets it free to concentrate on more advanced problems, and,
  in effect, increases the mental power of the
  race.}{\href{https://mathshistory.st-andrews.ac.uk/Biographies/Whitehead/}{Alfred
    North Whitehead} (1861-1947)}{.75}


\section{Historical Introduction}
\label{sec:hist-intr}
\markboth{\sc Differentials,  Differentiation, and the Derivative}{\sc Historical Introduction}

% \begin{wraptable}{r}{3.6in}
%   \begin{tabular}{ |c|c|c| }
      %       \hline
      %       %       \multicolumn{3}{|c|}{} \\
      %       \multicolumn{3}{|c|}{\bf{}General} \\
      %       \multicolumn{3}{|c|}{\bf{}Differentiation Rules} \\
      %       %       \multicolumn{3}{|c|}{} \\
      %       \hline
      %       %       && \\
      %       \small    Rule \#1&\small    The Constant Rule &\small If $a$ is a constant then \\
      %                       &\small     &\small $\d a = 0.$\\\hline
      %       %                       && \\
      %       \small    Rule \#2&\small    The Constant        &\small If $a$ is a constant and\\
      %                       &\small   Multiple Rule.       &\small $x$ is a variable then \\
      %                       &\small                        &\small$\d(ax) = a\d{x}.$\\\hline
      % %                       && \\
      %       \small    Rule \#3&\small    The Sum             &\small If $x$ and $y$ are variables then \\
      %                       &\small       rule.            &\small $\d(x+y) = \d x + \d y.$\\\hline{}
      %       %                       && \\
      %       \small    Rule \#4&\small    The Product         &\small If $x$ and $y$ are variables then\\
      %                       &\small     Rule.              &\small $\d(xy) = x\d y + y\d x.$\\\hline
      %       %                       && \\
      %       \small    Rule \#5&\small    The Power           &\small
                                                                 %                                                                  If $x$
                                                                 %                                                                  is a variable and\\
      %                       &\small      Rule.             &\small  $n$ is a rational number then\\
      %                       &\small                        &\small  $\d(x^n) = nx^{n-1}\d x$\\
      %       %                       && \\
      %       \hline{}
      %       %       && \\
      %       \small    Rule \#6&\small    The Quotient        &\small If $x$ and $y$ are variables then\\
      %                       &\small           Rule.        &\small  $\d\left(\frac{x}{y}\right) = \frac{y\d x
                                                               %                                                                - x\d y}{x^2}$\\
                                                               %                                                          %                                                                && \\
      %       \hline
      %       % %       && \\
      %       %       \small    Rule \#7&\small    The Chain        &\small If $y$ depends on
      %       %       c    $x,$ \\
      %       %       &\small    Rule.            &\small and $x$ depends on $t$ then\\
      %       %       &\small                     &\small  $\d y = \dfdx{y}{x}\dfdx{x}{t}\d t$\\[2mm]
      %       % %       && \\
      %       %       \hline
      %     \end{tabular}
      %       \end{wraptable}





\begin{wrapfigure}[11]{r}{3in}
  \vskip-5mm{}
  \captionsetup{labelformat=empty}
  \centerline{\includegraphics*[height=1.5in,width=3in]{../Figures/WoolsthorpeManor}}
  \caption{\textcolor{black}{\bf{}{Woolsthorpe Manor today}}}
  %\href{https://mathshistory.st-andrews.ac.uk/Diagrams/Woolsthorpe.jpeg}
  \label{fig:WoolsthorpeManor}
\end{wrapfigure}
During the Covid-19 pandemic of $2020$ the following story has become
very popular. In $1665$ an outbreak of Bubonic Plague (also called the
Black Death) closed Cambridge University.  One of the students, a
young man named
\href{https://mathshistory.st-andrews.ac.uk/Biographies/Newton/}{Isaac
  Newton}, decided to continue his studies on his own at his country
home in Woolsthorpe, Lincolnshire, UK. It was during this time that
Newton worked out the basic principles of Optics, his Law of Universal
Gravitation, and his version of the Calculus. He was under $25$ years
old at the time.

The basic facts of this story are essentially correct. The Plague did
close Cambridge University in $1665$, Newton was under $25$, he did
retire to Woolsthorpe to continue his studies, and he did investigate
Optics, Gravity, and Calculus during this time.

But telling Newton's story this way is misleading. It leaves the
impression that he did all of these things during his two years at
Woolsthorpe; sewed them up in a neat little package; bequeathed them
to the rest of humanity and then ascended directly into heaven as a
scientific demi-god.

\begin{wrapfigure}[18]{l}{2.1in}
  \captionsetup{labelformat=empty}
  \centerline{\href{https://mathshistory.st-andrews.ac.uk/Biographies/Newton/}{\includegraphics*[height=2.5in,width=2.1in]{../Figures/Newton}}}
  \caption{\centerline{Isaac Newton}\\\centerline{ 1642-1727}}
  \label{fig:Newton}
\end{wrapfigure}
This is not true\aside{To be sure Newton did a \emph{lot} during his
  retreat to Woolsthorpe, and he was a scientific demi-god. But if he
  ascended into heaven he did it in the usual way in $1727$ at the age
  of $84$.}. In fact,  he continued to work on optics after returning to
Cambridge, but did not publish his first results until $1672$. He
worked on his gravitational theory until at least $1687$ when he
published it in the first edition of his book \pubtitle{Philosophiae
  Naturalis Principia Mathematica} (Mathematical Principles of Natural
Philosopy\aside{It is usually just called \pubtitle{The Principia}.}).
The mathematics that he needed did not exist at the time
so he invented it.

When you're a genius you can do that.

Newton's story is often treated as inspirational but the fact that he
did so much, so young seems to us (the authors of this book) more
intimidating than inspiring. We are considerably older than $25$ and
taken together we haven't done anything as impressive as Newton did
in his youth.

But comparing ourselves to Newton is pointless. Newton was a true
genius.  We are merely ordinary people just like you. We didn't invent
Calculus, and you don't have to either. Newton did that for
us. Everyone since has merely had to learn it; a much less imposing
task.

Well, almost everyone.


\begin{wrapfigure}[17]{r}{2in}
  \vskip-7mm{}
  \captionsetup{labelformat=empty}
  \centerline{\href{https://mathshistory.st-andrews.ac.uk/Biographies/Leibniz/}{\includegraphics*[height=2.5in,width=2in]{../Figures/Leibniz}}}
  \caption{ \centerline{Gottfried Wilhelm Leibniz},\\
      \centerline{1646-1716}}
  \label{fig:Leibniz}
\end{wrapfigure}
About ten years after Newton entered Cambridge University, and quite
independently,
\href{https://mathshistory.st-andrews.ac.uk/Biographies/Leibniz}{Gottfried
  Wilhelm Leibniz} developed his rules for Calculus while employed as a
diplomat in Paris from $1672$ to $1674.$ He was $26$ years old at the
time.

Though their approaches are fundamentally different, the computational
techniques they worked out were essentially equivalent.  Newton's
approach was dynamic. He was thinking about objects in motion: the
moon, the planets, and falling apples. Leibniz was thinking about
geometric relationships, more like Fermat or Descartes.

As a matter of pure logic it doesn't really matter which viewpoint is
adopted. However, there are some situations where thinking dynamically
(like Newton) is more helpful, and others where thinking geometrically
(like Leibniz) is more helpful. Throughout this text we will switch
back and forth, adopting whichever approach seems to us more
advantageous for a given situation. You should too.

Newton was a private, almost secretive man. At first he didn't tell
anyone about his invention so it fell to Leibniz to publish the first
paper on Calculus.  In $1684$ he published a paper with the rather
imposing title
\begin{center}
  \begin{minipage}{.8\linewidth}
    \pubtitle{A New Method for Maxima and Minima, as Well as Tangents,
      Which is Impeded Neither by Fractional nor Irrational
      Quantities, and a Remarkable Type of Calculus for This.}
  \end{minipage}
\end{center}

  In this publication Leibniz displayed all of the General
Differentiation Rules seen in the table at the beginning of the next
section. He called his computational rules  ``\foreign{Calculus
  Differentialis}'' which translates loosely as ``rules for
differences.''  In English this is called \term{Differential
  Calculus}.

Each man created and used a notation that reflected his own viewpoint.
Newton's notation has (mostly) fallen away over the years. Leibniz'
notation is much more intuitive and has become standard in
mathematics.

Both men assumed that a varying quantity could be increased or
decreased by an infinitely small amount. Leibniz called this an
infinitesimal difference, or a \term{differential} and denoted it with
the letter ``d'' -- presumably for ``differential''.  So if $x$ is
some quantity, then its differential, $\dx{x}$, represents an
infinitely small increase or decrease in $x.$

To talk of infinitely small increments may sound strange to you. It
does to most people, because it is in fact a truly weird
idea. \href{https://mathshistory.st-andrews.ac.uk/Biographies/Bernoulli_Johann/}{Johann
  Bernoulli}\aside{Johann and his brother Jacob both learned Calculus
  from Leibniz, and both went on to mathematical greatness in their
  own right} described differentials this way:
``a quantity which is increased or decreased by an infinitely small
quantity is neither increased nor decreased\aside{From ``Great Moments
  in Mathematics After 1650'', Lecture 32}.''  You may judge for
yourself how helpful this is.

Before the invention of Calculus mathematicians had been using the
notion of an infinitesimal both successfully, and informally, for many
years. Galileo and his students used them. So did Archimedes.  The
idea was ``in the air.''  Late in his life Newton repudiated the idea
of differentials (at least, he appeared to) and replaced them with
what he called \emph{prime and ultimate ratios.}  Leibniz simply
continued the tradition of using differentials and invented a notation
to formalize it which we still use today. So for now we will be following
Leibniz' lead.

Although neither of them could answer the question, ``What is a
differential?''neither Leibniz nor Newton allowed this to hold them
back. They developed their computational schemes, verified that they
worked -- at least they worked on the problems they were interested in
-- and proceeded to use
them to attack other problems. That was enough for them. Both men
followed the example of their forebears and treated the notion of an
infinitely small increment as a ``convenient fiction;'' as something
that doesn't exist in reality but is useful as way to think about
certain problems.  We will adopt a similar attitude.


At first.

But this is mathematics, after all. We \emph{must} define our concepts
clearly and precisely. \emph{All} of them, not just the convenient
ones.


The question, ``What is a differential?''  cannot simply be swept
under the rug and ignored indefinitely. If we do not proceed
carefully, logically, from clearly stated first principles we can
never truly know if what we are doing is valid. Since everything we
do will depend in a fundamental way on differentials we \emph{must} resolve
this question or all of our work will rest on a foundation of shifting
sands and we will never really know whether it is valid or not. Indeed,
Newton, Leibniz, and generations of mathematicians since understood
this very well and worked hard to either give rigorous meaning to the
idea of a differential or to abandon it altogether in favor of some
other concept. Newton's idea of prime and ultimate ratios was probably
the first such attempt.

But it took some of the greatest minds in history nearly two hundred
years to resolve the question of differentials.

\begin{wrapfigure}[19]{l}{2in}
  \vskip-3mm{}
  \captionsetup{labelformat=empty}
  \centerline{\href{https://mathshistory.st-andrews.ac.uk/Biographies/Robinson/}{\includegraphics*[height=2.5in,width=2in]{../Figures/Robinson_4}}}
  \caption{\centerline{Abraham Robinson,}\\\centerline{1918-1974}}
  \label{fig:Robinson}
\end{wrapfigure}
Well, ``resolve'' might be too strong a word.  The theory of limits,
adopted in the late $19$th century, doesn't resolve the matter as much
as it simply avoids it altogether. You'll see what we mean when we
come to Chapter~\ref{chapt:whats-wrong-with}.

In the mid-twentieth century
\href{https://mathshistory.st-andrews.ac.uk/Biographies/Robinson}{Abraham
  Robinson} ($1918-1974$) took on the problem of differentials and was
finally able to put a solid, rigorous foundation underneath the use of
differentials in Calculus. Unfortunately, to do this Robinson required
some high level techniques from the field of Mathematical Logic which
are quite beyond the scope of this text so we will not be studying
Robinson's work. 

Besides, Calculus already had a rigorous foundation. The theory of limits
had been established nearly a hundred years before and it
provides an entirely adequate foundation for Calculus.  Unfortunately,
the theory of limits does very little to help anyone actually
\emph{use} Calculus. It simply provides rigor for the logical
foundations. On the other hand the differentials of Leibniz and
Robinson do provide aid in using Calculus. Newton, Leibniz and their
contemporaries did truly amazing things by treating differentials as
``convenient fictions.''  There is no reason we can't do the same.

So \underline{for} \underline{now}, we will follow the example of our
forbears and develop the ``Rules For Differences''
(\alert{Differential Calculus}) using Leibniz' differentials.  In
particular, we will not concern ourselves with the question, ``What is
a differential?''

There is precedent for this. The ancient Greek mathematician
\href{https://mathshistory.st-andrews.ac.uk/Biographies/Euclid/}{Euclid},
(325BC - 265BC)
defined geometric points with the phrase ``A point is that which has
no parts.'' But that is simply a clever way of admitting that he
didn't know what points are, despite the fact that they are enormously
useful.  What is a point, really, but a convenient fiction? And yet
the notion of a point is foundational for all of Geometry. After all,
what is a line but the relationship between two points? Moreover,
mathematical points are highly intuitive and can be visualized in the
mind's eye as long as we don't look too closely.

Differentials will play a similar role in Calculus.


%%% Local Variables: 
%%% mode: latex
%%% outline-minor-mode: t
%%% eval:(flyspell-mode)
%%% TeX-master: "DiffCalc"
%%% End: 
